\chapter{Analytic background}
\label{chap:analysis}

This chapter lists the general analytic facts the argument relies on, sorted
by how much Sobolev theory they need. Mathlib has Fourier analysis on the torus
(\texttt{UnitAddTorus}, \texttt{mFourierBasis}), absolutely continuous
functions on intervals, and Sobolev spaces $H^{s,p}(\R^n)$ defined through
tempered distributions. It has no Sobolev spaces on open subsets, no traces on
curves, no coarea formula, and no Sobolev embeddings into H\"older spaces.

\paragraph{Sobolev tiers.} Every node whose \emph{statement} involves a Sobolev
space uses one of the three definitions below, so the dependency graph shows
where Sobolev spaces enter.
\begin{enumerate}
  \item Definition~\ref{def:h1-interval}: $H^1$ on an interval, that is,
        absolutely continuous functions with square integrable derivative.
        Classical, and available in Mathlib in this form.
  \item Definition~\ref{def:sobolev-torus}: $H^s(\T^2)$ through Fourier
        coefficients. Cheap to define and use ourselves.
  \item Definition~\ref{def:sobolev-local}: Sobolev spaces on open sets,
        weighted spaces and traces. This is the real infrastructure gap.
\end{enumerate}
The main theorem's statement and the sector reduction use none of them.
The candidates for the initial \emph{assumptions} of the formalization are
the lemmas of Section~\ref{sec:local-sobolev}: they would be stated in a single
Lean file, each with a reference, and proved later or never.

\section{Function spaces}

\begin{definition}[$H^1$ on an interval]
\label{def:h1-interval}
For an interval $I\subset\R$, $H^1(I)$ is the space of functions that are
absolutely continuous on every compact subinterval, with $F,F'\in L^2(I)$;
$H^2(I)$ requires in addition $F'\in H^1(I)$. For bounded $I$ these functions
extend continuously to the endpoints, and $H^1_0(I)$ consists of those
vanishing there. The periodic version on a circle (an orbit of the flow) is
defined in the same way.
\end{definition}

\begin{definition}[Sobolev spaces on the torus]
\label{def:sobolev-torus}
\uses{def:torus}
For $s\ge0$, $H^s(\T^2)=\{g\in L^2:\sum_k(1+|k|^2)^s|\hat g_k|^2<\infty\}$.
Derivatives are the Fourier multipliers $ik_j$.
\end{definition}

\begin{definition}[Local Sobolev spaces and traces]
\label{def:sobolev-local}
\uses{def:sobolev-torus}
For an open $U\subset\T^2$ or $U\subset\R^2$ and an integer $k\ge0$, $H^k(U)$ is
the space of $L^2(U)$ functions whose weak derivatives up to order $k$ are in
$L^2(U)$, and $H^k_{\mathrm{loc}}(U)$ the space of functions in $H^k(V)$ for
every $V\Subset U$. Weighted versions require the weak derivatives to be in a
weighted $L^2$ space. Traces on curves are defined in Lemma~\ref{lem:traces}.
The fractional space $H^{1/2}(\R)$ is Mathlib's $H^{s,2}(\R)$; its norm is
equivalent to the one given by the Gagliardo seminorm.
\end{definition}

\section{Sobolev spaces on the torus}

\begin{lemma}[Elliptic regularity on the torus]
\label{lem:elliptic-torus}
\uses{def:sobolev-torus, def:inverse-laplacian}
$\Delta^{-1}:H^s\cap X\to H^{s+2}$ is bounded for $s\ge0$, and
$\partial_j:H^{s+1}\to H^s$. In particular $\psi=\Delta^{-1}\omega\in H^2$ for
$\omega\in X$.
\end{lemma}

\begin{proof}
Fourier multipliers.
\end{proof}

\begin{lemma}[Sobolev embeddings on the torus]
\label{lem:sobolev-embedding}
\uses{def:sobolev-torus}
$H^2(\T^2)\subset C^{0,\alpha}(\T^2)$ for every $\alpha<1$, and
$H^1(\T^2)\subset L^p(\T^2)$ for every $p<\infty$, with continuous inclusions.
\end{lemma}

\section{Local Sobolev theory (candidate assumptions)}
\label{sec:local-sobolev}

\begin{lemma}[Localization]
\label{lem:local-to-global}
\uses{def:sobolev-local}
If $u\in H^k_{\mathrm{loc}}(U)$ and $\chi$ is smooth with compact support in
$U$, then $\chi u\in H^k(\T^2)$ (extended by zero). Conversely the restriction
of an $H^k(\T^2)$ function to $U$ lies in $H^k(U)$, and weak derivatives agree
with the Fourier derivatives of Definition~\ref{def:sobolev-torus}.
\end{lemma}

\begin{lemma}[Local elliptic regularity]
\label{lem:elliptic-local}
\uses{def:sobolev-local}
If $u\in H^2$ near a point and $\Delta u\in H^1$ there, then $u\in H^3$ there.
Dyadic annular estimates of the form
$\|D^3u\|_{L^2(A_r)}\le C(\|\nabla\Delta u\|_{L^2(A_r^*)}
+r^{-1}\|\Delta u\|_{L^2(A_r^*)}+r^{-3}\|u\|_{L^2(A_r^*)})$ hold, with constants
independent of $r$.
\end{lemma}

\begin{lemma}[Traces]
\label{lem:traces}
\uses{def:sobolev-local}
Restriction to a smooth (or uniformly Lipschitz) curve maps $H^1$ to
$H^{1/2}$ and $H^2$ to $H^{3/2}$, continuously and locally uniformly in
families of such curves. In particular the traces of $\nabla\psi$ on the
regular level sets of $\Psi$ depend continuously on the level.
\end{lemma}

\section{Measure theory and operator theory}

\begin{lemma}[Coarea and a.e.\ restrictions]
\label{lem:coarea}
\uses{def:cells, def:h1-interval}
On a compact regular orbit annulus $\{\epsilon_1\le\Psi\le\epsilon_2\}$, for
$g\in L^1$,
$\int g\,dx\,dy=\int\!\big(\oint_{\{\Psi=c\}}g\,d\ell/\rho\big)dc$, so the
restriction of an $L^2$ function to almost every level is defined and square
integrable in flow time. If $g\in L^2$ and $L_0g\in L^2$, then on almost every
level the restriction is in $H^1$ of the orbit, with derivative $(L_0g)|$ in
flow time.
\end{lemma}

\begin{remark}
On a compact regular annulus the orbit coordinates (level, flow time) are a
smooth diffeomorphism onto a product, so Lemma~\ref{lem:coarea} reduces to
change of variables and Fubini; no general coarea formula is needed.
\end{remark}

\begin{lemma}[Self-adjointness of the transport operator]
\label{lem:H-selfadjoint}
\uses{def:dipole}
$H=-iL_0$ with domain $\{u\in L^2:L_0u\in L^2\}$ is self-adjoint on $L^2(\T^2)$,
and smooth functions form a core. It generates the unitary group
$u\mapsto u\circ\Phi_t$, where $\Phi_t$ is the (area preserving) flow of
$(\sin y,-\sin x)$.
\end{lemma}

\begin{remark}
Lemma~\ref{lem:H-selfadjoint} is the classical fact that a smooth divergence
free vector field on a compact manifold gives an essentially skew-adjoint
operator on smooth functions, whose closure is the maximal (distributional)
operator. Together with Stone's theorem this is the main functional-analytic
input to the virial identities. It needs no Sobolev spaces.
\end{remark}

\section{Elementary inequalities}

\begin{lemma}[One-dimensional Hardy inequalities]
\label{lem:hardy-1d}
\uses{def:h1-interval}
\begin{enumerate}
  \item If $F\in H^1(0,\eta)$ and $F(0)=0$, then
        $\int_0^\eta|F/x|^2\,dx\le4\int_0^\eta|F'|^2\,dx$.
  \item If $N\in H^2$ on an interval containing $x_0$ and $N(x_0)=0$, then
        $\|(N/(x-x_0))'\|_{L^2}\le C\|N''\|_{L^2}$; this follows from
        $N(x)/(x-x_0)=\int_0^1N'(x_0+t(x-x_0))\,dt$.
  \item (Angular Hardy) If $g\in H^1(0,\pi/2)$ vanishes at both endpoints, then
        $\int_0^{\pi/2}|g|^2/(\sin^2\theta\cos^2\theta)\,d\theta\le
        C\int_0^{\pi/2}|g'|^2\,d\theta$.
\end{enumerate}
\end{lemma}

\begin{lemma}[Critical Hardy inequality for odd functions]
\label{lem:odd-hardy}
\uses{def:sobolev-local}
If $a\in H^{1/2}(\R)$ is odd and compactly supported, then
$\int_0^\infty|a(u)|^2\,du/u\le C\|a\|_{H^{1/2}(\R)}^2$
(manuscript \texttt{(ee:odd-hardy)}).
\end{lemma}

\begin{proof}
The cross-half-line part of the Gagliardo seminorm contains
$\int_0^\infty\!\int_0^\infty|a(u)+a(v)|^2(u+v)^{-2}\,du\,dv$. The mixed term
is nonnegative since $(u+v)^{-2}=\int_0^\infty te^{-tu}e^{-tv}\,dt$, and the
two diagonal terms give $2\int_0^\infty|a(u)|^2\,du/u$.
\end{proof}

\begin{lemma}[Picone identity]
\label{lem:picone}
Let $G$ be a symmetric matrix field and $g>0$ smooth. Then pointwise
\[
  G\nabla z\cdot\nabla\bar z
  =g^2\,G\nabla(z/g)\cdot\nabla(\bar z/g)+G\nabla g\cdot\nabla(|z|^2/g).
\]
Consequently, after integration by parts on a domain $U$,
\[
  \int_U G\nabla z\cdot\nabla\bar z-V|z|^2
  =\int_U g^2G\nabla(z/g)\cdot\nabla(\bar z/g)
   +\int_U\frac{(-\Div(G\nabla g)-Vg)}{g}|z|^2
   +\int_{\partial U}\frac{\nu\cdot G\nabla g}{g}|z|^2 .
\]
\end{lemma}

\begin{proof}
Expand $\nabla(z/g)=(\nabla z)/g-z\nabla g/g^2$.
\end{proof}
